\documentclass[tikz,border=6pt]{standalone}
\usepackage{ctex}
\usepackage{amsmath}
\usetikzlibrary{calc, intersections}

\begin{document}
\begin{tikzpicture}[scale=1.2, thick]

% 坐标轴（浅灰色参考）
\draw[gray!35, thin, ->] (-2.8, 0) -- (4.2, 0) node[right, font=\footnotesize, gray] {$x$};
\draw[gray!35, thin, ->] (0, -2.4) -- (0, 3.2) node[above, font=\footnotesize, gray] {$y$};

% 圆 C1：圆心 O1(-0.8, 0)，半径 1.8
\def\Oax{-0.8} \def\Oay{0.0} \def\ra{1.8}
% 圆 C2：圆心 O2(1.6, 0)，半径 1.6
\def\Obx{1.6}  \def\Oby{0.0} \def\rb{1.6}

\draw[blue, thick, name path=c1] (\Oax, \Oay) circle (\ra);
\filldraw[blue] (\Oax, \Oay) circle (1.5pt);
\node[below, font=\small, blue] at (\Oax, \Oay) {$O_1$};

\draw[red, thick, name path=c2] (\Obx, \Oby) circle (\rb);
\filldraw[red] (\Obx, \Oby) circle (1.5pt);
\node[below, font=\small, red] at (\Obx+0.05, \Oby) {$O_2$};

% 公共弦所在直线（由两圆方程相减得到）
% C1: (x+0.8)^2 + y^2 = 3.24  → x^2+y^2+1.6x-2.6=0
% C2: (x-1.6)^2 + y^2 = 2.56  → x^2+y^2-3.2x-0.16=0
% 相减：4.8x - 2.44 = 0  → x = 2.44/4.8 ≈ 0.5083
% 公共弦是竖直线 x ≈ 0.508，但为了图示更清楚采用略微倾斜
% 实际计算：设 C1: x^2+y^2+1.6x = 2.6, C2: x^2+y^2-3.2x = 0.16
% 相减: 4.8x = 2.44 → x = 61/120 ≈ 0.5083
% 公共弦是 x = 61/120，与 C1 求交点：
% (61/120 + 0.8)^2 + y^2 = 3.24
% (61/120 + 96/120)^2 + y^2 = 3.24
% (157/120)^2 + y^2 = 3.24
% y^2 = 3.24 - 1.7057 = 1.5343 → y ≈ ±1.2387
\def\xchord{0.508}
\def\yA{1.239}
\def\yB{-1.239}

% 公共弦延长线（虚线）
\draw[green!50!black, dashed, thin] (\xchord, -2.0) -- (\xchord, 2.8)
  node[above, font=\footnotesize, green!50!black] {公共弦所在直线};

% 公共弦 AB（实线，加粗）
\draw[green!40!black, very thick] (\xchord, \yB) -- (\xchord, \yA);

% 交点 A, B
\filldraw[black] (\xchord, \yA) circle (2pt);
\node[right, font=\small] at (\xchord+0.05, \yA+0.05) {$A$};

\filldraw[black] (\xchord, \yB) circle (2pt);
\node[right, font=\small] at (\xchord+0.05, \yB-0.1) {$B$};

% 圆心连线
\draw[gray, thin, dashed] (\Oax, \Oay) -- (\Obx, \Oby);
\node[above, font=\footnotesize, gray] at (0.4, 0.1) {$d$};

% 公共弦长标注
\draw[gray!70, thin, <->] (\xchord-0.35, \yB) -- (\xchord-0.35, \yA)
  node[midway, left, font=\footnotesize, black] {$|AB|$};

% 说明文字
\node[font=\footnotesize, align=left, anchor=north west] at (-2.7, -1.6)
  {两圆方程相减消去 $x^2+y^2$\\得公共弦所在直线方程};

\end{tikzpicture}
\end{document}
